\documentclass[a4paper]{article}
%\usepackage{simplemargins}

%\usepackage[square]{natbib}
\usepackage{amsmath}
\usepackage{amsfonts}
\usepackage{amssymb}
\usepackage{graphicx}

\begin{document}
\pagenumbering{gobble}

\Large

\title{Joint Inversion of Magnetotelluric and P-Wave Travel Time Data at Mount St. Helens Volcano Using a Variation of Information Constraint}

% \hspace{10pt}

% % Author names and affiliations
% \large
% Katharina Maetschke$^1$, Max Moorkamp$^2$, Graham Hill$^3$, Lindsey Heagy$^1$\\
% \hspace{10pt}

% \small  
% $^1$) University of British Columbia, Department of Earth, Ocean and Atmospheric Sciences, Vancouver, BC, Canada\\
% $^2$) Technical University Berlin, Faculty VI, Applied Geophysics, Berlin, Germany\\
% $^3$) Czech Academy of Sciences, Institute of Geophysics, Prague, Czech Republic\\

% \end{center}

We present a first case study applying a Variation of Information (VI) coupling constraint in the joint inversion of magnetotelluric (MT) and P-wave travel time data in a volcanological setting. Our study utilizes previously published data sets from Mount St. Helens, USA, which were individually inverted. By combining two geophysical methods that are both sensitive to the presence of melt and hydrothermal fluids, we demonstrate that the VI coupling effectively reduces ambiguities in the interpretation of anomalies. Additionally, it permits obtaining parameter relationships without any a priori assumptions about the connection between conductivity and seismic velocity. Clustering analyses of these derived parameter relationships provide new insight into the spatial distribution of features characterized by specific conductivity and P-wave velocity relationships, which enhances interpretability. Our approach highlights the potential of VI-based coupling in enhancing the information gain and reducing interpretational ambiguities when using joint inversions in complex volcanic environments.
\end{document}